% chapters/chap3.tex

\chapter{使い分けのまとめ}

\section{どちらを使うか}

\begin{table}[h]
  \centering
  \begin{tabular}{lll}
    \hline
    用途 & 推奨コマンド & 理由 \\
    \hline
    マクロ・設定の分割    & \texttt{\textbackslash input}   & プリアンブルで使えるため \\
    節（section）の分割   & \texttt{\textbackslash input}   & 改ページなしで連続するため \\
    章（chapter）の分割   & \texttt{\textbackslash include} & 章ごとに改ページが自然なため \\
    \hline
  \end{tabular}
  \caption{\texttt{\textbackslash input} と \texttt{\textbackslash include} の使い分け}
\end{table}

\section{\texttt{\textbackslash includeonly} の使い方}

\begin{enumerate}
  \item 通常は \texttt{\textbackslash includeonly} をコメントアウトして全章コンパイル
  \item 特定の章を編集するときだけ \texttt{\textbackslash includeonly} を有効にする
  \item 作業が終わったら再度コメントアウトして全章コンパイル
\end{enumerate}
